\section{Questao 6 - Monte Carlo Tree Search (MCTS)}

\subsection{Escopo e modelagem usada}
Para a Q6, foi usada uma modelagem simplificada e autoral do estado de decisao na raiz do Connect-4:
\begin{itemize}[leftmargin=1.5em]
  \item Acoes possiveis na raiz: $\{C1,C2,C3,C4\}$;
  \item cada acao gera um filho imediato na arvore de busca;
  \item cada simulacao retorna resultado binario para MAX (vermelho): vitoria = 1, derrota = 0.
\end{itemize}

Representacao em Python (\texttt{codes/q6\_mcts.py}):
\begin{itemize}[leftmargin=1.5em]
  \item no armazena \texttt{visits} $=N(s)$ e \texttt{wins} $=W(s)$;
  \item estrutura de arvore por ponteiro \texttt{parent} + dicionario \texttt{children};
  \item implementacao 100\% autoral, sem bibliotecas externas de jogos/MCTS.
\end{itemize}

\subsection{Explicacao do MCTS (itens teoricos do enunciado)}
\begin{itemize}[leftmargin=1.5em]
  \item \textbf{Selecao:} desce pela arvore escolhendo filho com maior UCT.
  \item \textbf{Expansao:} ao atingir no com filhos ainda nao criados, adiciona-se um novo filho.
  \item \textbf{Simulacao (rollout):} a partir do no expandido, simula-se uma partida ate termino (aqui, modelo simplificado de resultado).
  \item \textbf{Retropropagacao:} atualiza-se $N(s)$ e $W(s)$ do no atual ate a raiz.
\end{itemize}

Metricas pedidas:
\begin{itemize}[leftmargin=1.5em]
  \item $N(s)$: numero de visitas do estado.
  \item $W(s)$: numero de vitorias acumuladas para MAX.
  \item Exploitation: $\frac{W(s)}{N(s)}$ (qualidade media observada).
  \item Exploration: termo que prioriza nos pouco visitados.
\end{itemize}

\subsection{Formula UCT usada}
\[
\mathrm{UCT}(s_i)=\frac{W(s_i)}{N(s_i)} + C\sqrt{\frac{\ln N(\text{pai})}{N(s_i)}}
\]
com $C=1.4$ no experimento principal.

\subsection{Execucao detalhada (10 iteracoes)}
\begingroup
\footnotesize
\setlength{\tabcolsep}{4pt}
\begin{longtable}{c p{2.5cm} p{3.0cm} c c}
\toprule
Iter. & Caminho & No expandido/reusado & Rollout & $(N,W)$ no no escolhido \\
\midrule
1 & raiz $\rightarrow C1$ & $C1$ & vitoria & $(1,1)$ \\
2 & raiz $\rightarrow C2$ & $C2$ & vitoria & $(1,1)$ \\
3 & raiz $\rightarrow C3$ & $C3$ & vitoria & $(1,1)$ \\
4 & raiz $\rightarrow C4$ & $C4$ & derrota & $(1,0)$ \\
5 & raiz $\rightarrow C1$ & (reuso) $C1$ & vitoria & $(2,2)$ \\
6 & raiz $\rightarrow C2$ & (reuso) $C2$ & vitoria & $(2,2)$ \\
7 & raiz $\rightarrow C3$ & (reuso) $C3$ & derrota & $(2,1)$ \\
8 & raiz $\rightarrow C1$ & (reuso) $C1$ & vitoria & $(3,3)$ \\
9 & raiz $\rightarrow C2$ & (reuso) $C2$ & derrota & $(3,2)$ \\
10 & raiz $\rightarrow C1$ & (reuso) $C1$ & derrota & $(4,3)$ \\
\bottomrule
\end{longtable}
\endgroup

\subsection{Arvore parcial apos as 10 iteracoes}
\begin{figure}[H]
\centering
\begin{tikzpicture}[
  level distance=20mm,
  sibling distance=22mm,
  rootnode/.style={draw, rounded corners, minimum width=24mm, minimum height=10mm, align=center, fill=gray!10},
  bestnode/.style={draw, circle, minimum size=12mm, align=center, thick, blue!70},
  nodeok/.style={draw, circle, minimum size=12mm, align=center, fill=gray!10},
  edge from parent/.style={-Latex}
]
\node[rootnode] (r) {raiz\\$N=10,\ W=6$}
  child { node[bestnode] {$C1$\\\scriptsize $N=4,W=3$} }
  child { node[nodeok] {$C2$\\\scriptsize $N=3,W=2$} }
  child { node[nodeok] {$C3$\\\scriptsize $N=2,W=1$} }
  child { node[nodeok] {$C4$\\\scriptsize $N=1,W=0$} };
\end{tikzpicture}
\caption{Arvore parcial do MCTS apos 10 iteracoes (azul = jogada mais visitada)}
\end{figure}

\subsection{Tabela UCT apos 10 iteracoes}
Para $N(\text{pai})=10$ e $C=1.4$:
\begingroup
\small
\begin{table}[H]
\centering
\begin{tabular}{c c c c c c}
\toprule
Jogada & $N$ & $W$ & $\frac{W}{N}$ & $1.4\sqrt{\frac{\ln 10}{N}}$ & UCT \\
\midrule
$C1$ & 4 & 3 & 0.7500 & 1.0622 & 1.8122 \\
$C2$ & 3 & 2 & 0.6667 & 1.2265 & 1.8932 \\
$C3$ & 2 & 1 & 0.5000 & 1.5022 & 2.0022 \\
$C4$ & 1 & 0 & 0.0000 & 2.1244 & 2.1244 \\
\bottomrule
\end{tabular}
\caption{Calculo de UCT por filho da raiz}
\end{table}
\endgroup

\textbf{Leitura da tabela:}
\begin{itemize}[leftmargin=1.5em]
  \item maior \textit{exploitation} (\(W/N\)): $C1$;
  \item maior UCT no instante: $C4$, por forte termo de exploracao (poucas visitas);
  \item decisao final de jogada no codigo: por maior numero de visitas, resultando em \textbf{$C1$}.
\end{itemize}

\subsection{Sensibilidade ao parametro C (rollout semi-guloso)}
\begingroup
\footnotesize
\setlength{\tabcolsep}{4pt}
\begin{table}[H]
\centering
\begin{tabular}{c p{2.8cm} c c c}
\toprule
$C$ & Movimento mais estavel & Estabilidade & Diversidade media & Winrate medio \\
\midrule
0.1 & $C1$ & 57.5\% & 0.42 & 0.696 \\
1.4 & $C1$ & 75.0\% & 0.94 & 0.641 \\
3.0 & $C1$ & 72.5\% & 0.99 & 0.627 \\
\bottomrule
\end{tabular}
\caption{Impacto do parametro de exploracao $C$}
\end{table}
\endgroup

\textbf{Interpretacao:}
\begin{itemize}[leftmargin=1.5em]
  \item $C=0.1$: baixa exploracao, menor diversidade e tendencia mais gulosa;
  \item $C=1.4$: melhor equilibrio entre explorar e consolidar decisao;
  \item $C=3.0$: exploracao muito alta, aumentando diversidade, mas com queda de qualidade media.
\end{itemize}

\subsection{Rollout aleatorio vs semi-guloso (C=1.4)}
\begingroup
\footnotesize
\setlength{\tabcolsep}{4pt}
\begin{table}[H]
\centering
\begin{tabular}{p{2.3cm} p{2.6cm} c c c}
\toprule
Politica de rollout & Movimento mais estavel & Estabilidade & Diversidade media & Winrate medio \\
\midrule
Aleatorio & $C1$ & 70.0\% & 0.94 & 0.551 \\
Semi-guloso & $C1$ & 75.0\% & 0.94 & 0.641 \\
\bottomrule
\end{tabular}
\caption{Comparacao entre politicas de simulacao}
\end{table}
\endgroup

\textbf{Conclusao desta comparacao:}
\begin{itemize}[leftmargin=1.5em]
  \item semi-guloso foi melhor em qualidade (maior winrate medio);
  \item semi-guloso tambem foi ligeiramente mais estavel na escolha final;
  \item mantendo mesma diversidade media neste experimento.
\end{itemize}

\subsection{Efeito do numero de iteracoes}
Experimento adicional com $C=1.4$:
\begingroup
\footnotesize
\setlength{\tabcolsep}{4pt}
\begin{table}[H]
\centering
\begin{tabular}{c p{2.3cm} c c c}
\toprule
Iteracoes & Rollout & Movimento estavel & Estabilidade & Winrate medio \\
\midrule
10 & semi-guloso & $C1$ & 65.0\% & 0.613 \\
20 & semi-guloso & $C1$ & 62.5\% & 0.611 \\
40 & semi-guloso & $C1$ & 75.0\% & 0.641 \\
80 & semi-guloso & $C1$ & 72.5\% & 0.631 \\
120 & semi-guloso & $C1$ & 82.5\% & 0.651 \\
\midrule
10 & aleatorio & $C1$ & 62.5\% & 0.525 \\
20 & aleatorio & $C1$ & 55.0\% & 0.536 \\
40 & aleatorio & $C1$ & 70.0\% & 0.551 \\
80 & aleatorio & $C1$ & 67.5\% & 0.552 \\
120 & aleatorio & $C1$ & 72.5\% & 0.562 \\
\bottomrule
\end{tabular}
\caption{Impacto do numero de iteracoes no MCTS}
\end{table}
\endgroup

Em geral, mais iteracoes aumentam estabilidade da decisao e qualidade media, com ganho mais claro no rollout semi-guloso.

\subsection{Comparacao final pedida na Q6}
\begin{itemize}[leftmargin=1.5em]
  \item \textbf{(a) Rollout aleatorio vs semi-guloso:} semi-guloso venceu em qualidade e estabilidade.
  \item \textbf{(b) Diferentes valores de C:} $C$ baixo explora pouco; $C$ alto explora bastante; $C=1.4$ foi o melhor compromisso observado.
  \item \textbf{(c) Numero de iteracoes:} aumentar iteracoes melhorou consolidacao da jogada escolhida.
  \item \textbf{(d) Estabilidade das decisoes:} maior com semi-guloso e maior orcamento de simulacoes.
\end{itemize}

\subsection{Resposta objetiva da jogada sugerida}
No cenario principal (10 iteracoes, $C=1.4$, rollout semi-guloso), a jogada escolhida pelo criterio de maior visita foi:
\[
\boxed{C1}
\]
